% !TeX program = pdflatex
\documentclass[12pt,a4paper]{article}

\newcommand{\notelanguage}{finnish}
% Shared preamble for course notes and exercise sets.

\usepackage[T1]{fontenc}
\usepackage{lmodern}
\usepackage[english]{babel}

\usepackage[a4paper,margin=30mm]{geometry}
\usepackage{microtype}
\usepackage{graphicx}
\usepackage{enumitem}

\usepackage{amsmath}
\usepackage{amssymb}
\usepackage{amsthm}
\usepackage{mathtools}
\usepackage{aliascnt}

\usepackage[hidelinks,pdfusetitle]{hyperref}

\numberwithin{equation}{section}
\setlist{itemsep=0.25em,topsep=0.5em}

\theoremstyle{plain}
\newtheorem{theorem}{Theorem}[section]
\newaliascnt{lemma}{theorem}
\newtheorem{lemma}[lemma]{Lemma}
\aliascntresetthe{lemma}
\newaliascnt{proposition}{theorem}
\newtheorem{proposition}[proposition]{Proposition}
\aliascntresetthe{proposition}
\newaliascnt{corollary}{theorem}
\newtheorem{corollary}[corollary]{Corollary}
\aliascntresetthe{corollary}

\theoremstyle{definition}
\newaliascnt{definition}{theorem}
\newtheorem{definition}[definition]{Definition}
\aliascntresetthe{definition}
\newaliascnt{example}{theorem}
\newtheorem{example}[example]{Example}
\aliascntresetthe{example}
\newaliascnt{problem}{theorem}
\newtheorem{problem}[problem]{Problem}
\aliascntresetthe{problem}

\theoremstyle{remark}
\newaliascnt{remark}{theorem}
\newtheorem{remark}[remark]{Remark}
\aliascntresetthe{remark}

\usepackage[nameinlink,noabbrev]{cleveref}

\crefname{theorem}{theorem}{theorems}
\crefname{lemma}{lemma}{lemmas}
\crefname{proposition}{proposition}{propositions}
\crefname{corollary}{corollary}{corollaries}
\crefname{definition}{definition}{definitions}
\crefname{example}{example}{examples}
\crefname{problem}{problem}{problems}
\crefname{remark}{remark}{remarks}

\DeclareMathOperator{\im}{im}
\DeclareMathOperator{\rank}{rank}
\DeclarePairedDelimiter{\abs}{\lvert}{\rvert}
\DeclarePairedDelimiter{\norm}{\lVert}{\rVert}
\DeclarePairedDelimiter{\set}{\lbrace}{\rbrace}

\usepackage{tikz}

\title{Analyysi I\\Harjoituksen 2 ratkaisut}
\author{}
\date{}

\begin{document}

\maketitle

\noindent

\begin{exercise}
  Olkoon \((x_n)\) lukujono, \(x_n = \frac{2n - 4}{n + 2}\). Osoita suoraan
  määritelmästä, että \(\lim_{n \to \infty} x_n = 2\).
\end{exercise}

\begin{solution}
  Olkoon \(\epsilon > 0\). Valitaan \(N \in \NN\) jolla \(N >
  \frac{8}{\epsilon}\). Kun \(n > N\), niin
  \[
    \abs{x_n - 2} = \frac{8}{n + 2} < \frac{8}{n} < \frac{8}{N} < \epsilon.
  \]
  Täten \(\lim_{n \to \infty} x_n = 2\).
\end{solution}

\begin{exercise}
  Olkoon \((x_n)\) lukujono, \(x_n = \frac{4n}{n - 3}\). Osoita suoraan
  määritelmästä, että \(\lim_{n \to \infty} x_n = 4\).
\end{exercise}

\begin{solution}
  Olkoon \(\epsilon > 0\). Valitaan \(N \in \NN\) jolla \(N >
  \maxset{6, \frac{24}{\epsilon}}\). Kun \(n > N\), niin
  \[
    \abs{x_n - 4} = \frac{12}{n - 3} \leq \frac{24}{n} < \frac{24}{N}
    < \epsilon.
  \]
  Täten \(\lim_{n \to \infty} x_n = 4\).
\end{solution}

\begin{exercise}
  Olkoon \((y_n)\) lukujono, \(y_n = (-1)^{n + 1}n\). Tutki määritelmän
  avulla suppenevatko jono \((y_n)\) ja alijono \((y_{2n})\).
\end{exercise}

\begin{solution}
  Tehdään vastaoletus: Oletetaan, että \(\lim_{n \to \infty} y_n = a\).
  On olemassa \(N \in \NN\) siten, että kaikilla \(n > N\) pätee
  \[
    \abs{y_n - a} < 1.
  \]
  Valitaan parillinen \(n \in \NN\) niin suuri, että \(n > \maxset{N,
  \abs{a} + 1}\).
  Koska \(n\) on parillinen, \(y_n = -n\), joten kolmioepäyhtälön nojalla,
  \[
    \abs{y_n - a} = \abs{-n - a} \geq n - \abs{a} > 1.
  \]
  Koska \(n > N\), oletuksen nojalla pitäisi olla \(\abs{y_{n} - a}\) < 1.
  Olemme saaneet ristiriidan, joten \((y_n)\) ei suppene.

  Tutkitaan seuraavaksi alijonoa \((y_{2n})\). Tehdään jälleen vastaoletus:
  Oletetaan, että \(\lim_{n \to \infty} y_{2n} = a\). On olemassa
  \(N \in \NN\) siten, että kaikilla \(n > N\) pätee
  \[
    \abs{y_{2n} - a} < 1.
  \]
  Valitaan \(n \in \NN\) niin suuri, että \(n > N + \abs{a} + 1\).
  Koska \(y_{2n} = -2n\), kolmioepäyhtälön nojalla
  \[
    \abs{y_{2n} - a} = \abs{-2n - a} \geq 2n - \abs{a} > 1.
  \]
  Koska \(n > N\), oletuksen nojalla pitäisi olla \(\abs{y_{2n} - a} < 1\).
  Olemme saaneet ristiriidan, joten \((y_{2n})\) ei suppene.
\end{solution}

\begin{exercise}
  Olkoon \((x_n)\) lukujono, joka suppenee kohti reaalilukua \(a\). Olkoon
  \(m \in \RR\), jolle \(x_n < m\) kaikilla \(n \in \NN\). Piirrä
  kuva tilanteesta. Osoita, että \(a \leq m\). Onko välttämättä
  \(a < m\)?
\end{exercise}

\begin{solution}
  Kuvassa vasemmalla raja-arvo on aidosti ylärajaa pienempi. Oikealla
  lukujonon termit lähestyvät ylärajaa alhaalta, joten raja-arvo voi myös olla
  yhtä suuri kuin yläraja.
  \begin{center}
    \begin{tikzpicture}[x=0.75cm,y=1.15cm, every node/.style={font=\small}]
      \begin{scope}
        \draw[->] (0,0) -- (6.8,0) node[right] {$n$};
        \draw[->] (0,0) -- (0,4.1);
        \draw[dashed] (0,3.45) -- (6.4,3.45) node[right] {$m$};
        \draw[dotted] (0,2.2) -- (6.4,2.2) node[right] {$a$};
        \foreach \x/\y in {1/0.9,2/2.8,3/1.65,4/2.5,5/1.95,6/2.32}
        \fill (\x,\y) circle (1.7pt);
        \node at (3.2,-0.55) {$a<m$};
      \end{scope}

      \begin{scope}[xshift=8.2cm]
        \draw[->] (0,0) -- (6.8,0) node[right] {$n$};
        \draw[->] (0,0) -- (0,4.1);
        \draw[dashed] (0,3.45) -- (6.4,3.45) node[right] {$a=m$};
        \foreach \x/\y in {1/1.1,2/2.05,3/2.55,4/2.85,5/3.05,6/3.18}
        \fill (\x,\y) circle (1.7pt);
        \node at (3.2,-0.55) {$a=m$};
      \end{scope}
    \end{tikzpicture}
  \end{center}

  Tehdään vastaoletus: Oletetaan \(a > m\). Olkoon \(\epsilon =
  \frac{a - m}{2} > 0\). Koska \((x_n)\) suppenee, on olemassa \(N \in \NN\),
  jolla
  \[
    \abs{x_n - a} < \epsilon,
  \]
  kun \(n > N\). Saadaan
  \[
    x_n > a - \epsilon = \frac{2a - (a - m)}{2} = \frac{a + m}{2} > m.
  \]
  Tämä on ristiriidassa oletuksen \(x_n < m\). On siis oltava
  \(a \leq m\).

  Osoitetaan vielä että \(a < m\) \emph{ei} pidä välttämättä paikkaansa.
  Olkoon \((x_n)\) lukujono, \(x_n = m - \frac{1}{n}\). Tällöin
  \(x_n < m\) kaikilla \(n \in \NN\). Näytetään
  \[
    \lim_{n \to \infty} x_n = m.
  \]
  Olkoon \(\epsilon > 0\). Valitaan \(N \in \NN\) siten, että
  \(N > \frac{1}{\epsilon}\). Kun \(n > N\), niin
  \[
    \abs{x_n - m} = \frac{1}{n} <
    \frac{1}{\bigl(\frac{1}{\epsilon}\bigr)} = \epsilon.
  \]
  Täten \(a = \lim_{n \to \infty} x_n = m\).
\end{solution}

\begin{theorem}\label{thm:convergent-sequences-are-bounded}
  Kaikki suppenevat lukujonot on rajoitettuja.
\end{theorem}

\begin{proof}
  Olkoon \((x_n)\) lukujono joka suppenee kohti reaalilukua \(a\). On
  olemassa \(N \in \NN\), jolle
  \[
    \abs{x_n - a} < 1,
  \]
  kun \(n > N\). Tästä saadaan yläraja
  \[
    \abs{x_n} < \abs{a} + 1,
  \]
  kun \(n > N\). Valitaan
  \[
    M = \maxset{\abs{x_0}, \cdots, \abs{x_N}, \abs{a} + 1}.
  \]
  Nyt epäyhtälö \(\abs{x_n} \leq M\) on voimassa kaikilla \(n \in \NN\),
  eli \((x_n)\) on rajoutettu.
\end{proof}

\begin{exercise}
  Olkoon \((x_n)\) lukujono, joka suppenee kohti reaalilukua \(a\), ja
  olkoon \((y_n)\) lukujono, joka suppenee kohti reaalilukua \(b\).
  Osoita, että lukujono \((x_ny_n)\) suppenee kohti lukua \(ab\).
\end{exercise}

\begin{solution}
  Olkoon \(\epsilon > 0\). Koska \((y_n)\) suppenee, on Lauseen
  \ref{thm:convergent-sequences-are-bounded} nojalla olemassa luku
  \(M > 0\), jolla \(\abs{y_n} < M\) kaikilla \(n \in \NN\). Valitaan
  \(N_1 \in \NN\) siten, että kun \(n > N_1\), niin
  \[
    \abs{x_n - a} < \frac{\epsilon}{2M}.
  \]
  Valitaan \(N_2 \in \NN\) siten, että kun \(n > N_2\), niin
  \[
    \abs{y_n - b} < \frac{\epsilon}{2(\abs{a} + 1)}.
  \]
  Valitaan \(N = \maxset{N_1, N_2}\). Kolmioepäyhtälön nojalla,
  kun \(n > N\), saadaan
  \[
    \begin{aligned}
      \abs{x_ny_n - ab} &\leq \abs{x_ny_n - ay_n} + \abs{ay_n - ab} \\
      &= \abs{y_n}\abs{x_n - a} + \abs{a}\abs{y_n - b} \\
      &< M\frac{\epsilon}{2M} + \abs{a}\frac{\epsilon}{2(\abs{a} + 1)} \\
      &< \frac{\epsilon}{2} + \frac{\epsilon}{2} \\
      &= \epsilon.
    \end{aligned}
  \]
  Täten \(\lim_{n \to \infty} x_ny_n = ab\).
\end{solution}

\begin{exercise}[Suppiloperiaate eli kuristusperiaate]
  Olkoon \((x_n)\), \((y_n)\) ja \((z_n)\) sellaisia lukujonoja, että
  \(x_n \leq y_n \leq z_n\) kaikilla \(n \in \NN\). Oletetaan, että
  lukujonot \((x_n)\) ja \((z_n)\) suppenevat ja
  \(\lim_{n \to \infty} x_n = a = \lim_{n \to \infty} z_n\).
  Osoita, että lukujono \((y_n)\) suppenee ja \(\lim_{n \to \infty} y_n = a\).
\end{exercise}

\begin{solution}
  Olkoon \(\epsilon > 0\). On olemassa \(N_1 \in \NN\), jolla
  \[
    \abs{x_n - a} < \epsilon \quad \text{kun} \quad n > N_1,
  \]
  sekä \(N_2 \in \NN\), jolla
  \[
    \abs{z_n - a} < \epsilon \quad \text{kun} \quad n > N_2.
  \]
  Valitaan \(N \in \NN\) siten, että \(N = \maxset{N_1, N_2}\). Kun
  \(n > N\), saadaan
  \[
    a - \epsilon < x_n \leq y_n \leq z_n < a + \epsilon,
  \]
  eli
  \[
    \abs{y_n - a} < \epsilon.
  \]
  Täten \(\lim_{n \to \infty} y_n = a\).
\end{solution}
\end{document}
